\subsection{The actual nonsplit block extension}
\label{subsec:sr-four-actual-block-extension}

We supply the additional producer used to pass from the
characteristic-\(\ell\) support calculation to finite weight zero
and the no-full-block conclusion.

\begin{srproposition}\label{prop:sr-four-actual-block-extension}
Let \(r\) be an actual integral Case II block-diagonal point
whose \(0,2\) block is irreducible over the fraction field.
There is an actual integral symplectic block extension \(r'\)
with the same characteristic polynomial at every group
element, the same residual representation, and the same
ordered inertial characters and ramification conditions.
Its other rank-two block is irreducible and the extension
does not split.
This holds both in mixed characteristic and for a
characteristic-\(\ell\) complete DVR model.
\end{srproposition}
\begin{proof}
Work over its complete DVR model \(\mathcal A\), with
uniformizer \(\pi\) and finite residue field \(k\).
After grouping indices as \((0,2),(1,3)\), let \(M\) be the
free rank-three \(\mathcal A\)-module of symplectic Lie
matrices in the lower off-diagonal block.
It is the symmetric-square block radical:
\(XY=0\) for \(X,Y\in M\), and \(X^\#=-X\).
Conjugation by \(r\) preserves \(M\).

At complex conjugation \(c\), the actual block matrix is
exactly the prescribed block scalar
\(\operatorname{diag}(-1,1,-1,1)\).
Indeed its square is \(1\), its reductions on the two blocks
are respectively \(-1,1\), and the opposite factor of
\(X^2-1\) is invertible on each block, since \(2\) is a unit.
Consequently \(\operatorname{Ad}r(c)=-1\) on \(M\).

The actual local separator \(\sigma_\ell\) satisfies
\[
 \bar\varepsilon(\sigma_\ell)=1,\qquad
 \operatorname{Ad}r(\sigma_\ell)-1
                         \text{ invertible on }M.         \tag{BE1}
\]
Its three residual eigencharacters are the odd roots
\(-\alpha,\beta,-\theta\); all are nontrivial at this
separator by the arithmetic data.
The full cyclotomic value is \(1\) here, as in the actual
chosen decomposition element.
In particular (BE1) also kills the invariants of
\(M^\vee(1)\), globally and locally.

Let \(u\) be the actual ordinary flag matrix and put
\[
 P=\{X\in M:u^{-1}Xu\text{ strictly upper triangular}\}.
\]
This is a saturated, stable free rank-one submodule.
We verify the rank assertion.
The two block characteristic polynomials of
\(r(\sigma_\ell)\) have disjoint residual roots, by the
cross-block separator. Their resultant is a unit.
The block projector is therefore a polynomial in
\(r(\sigma_\ell)\) over \(\mathcal A\), by Bezout and
Cayley--Hamilton. As \(r(\sigma_\ell)\) preserves the
ordinary flag, the projector does too.
The ordinary flag consequently decomposes into its two
block flags, with successive block labels \(0,1,2,3\).
A block-compatible symplectic basis identifies \(P\) with
the single upper-root line \(E_{12}\).
This also proves stability and saturation; alternatively
saturation follows directly from its defining zero-entry
equations. Thus \(M/P\) is free.

We construct a nonzero integral cocycle rather than
postulating an extension class.
For \(n\ge1\), set \(M_n=M/\pi^nM\), \(P_n=P/\pi^nP\).
At \(\ell\) impose the local condition
\[
 L_\ell=\operatorname{im}
           \bigl(H^1(\mathbb Q_\ell,P_n)
                         \longrightarrow H^1(\mathbb Q_\ell,M_n)\bigr).
\]
At every other finite prime impose the unramified condition.
The invariants of \(P_n,(M/P)_n\), and \(P_n^\vee(1)\)
vanish at \(\ell\), by (BE1): invertibility there is
preserved on a stable saturated submodule and its quotient.
Local duality and the Euler characteristic formula give
\[
 H^2(\mathbb Q_\ell,P_n)=0,\qquad
 |H^1(\mathbb Q_\ell,P_n)|=|P_n|,\qquad
                         |L_\ell|=|P_n|.                 \tag{BE2}
\]
The last map is injective because \(H^0((M/P)_n)=0\).

Take a finite global ramification set containing that of
\(M_n\), the original level and \(\ell\); it may be fixed
for all \(n\) from the original DVR representation.
At a prime different from \(\ell\),
\[
                      |H^1_{\rm unr}(\mathbb Q_p,M_n)|
                                =|H^0(\mathbb Q_p,M_n)|.
\]
Indeed the unramified group is the cokernel of Frobenius
minus one on the finite module \(M_n^{I_p}\), whose
cokernel and kernel have equal cardinality.
The real local factor is one, since complex conjugation
acts as \(-1\) and \(\ell\ne2\).
The global invariants and dual invariants are zero by
(BE1).
The exact Greenberg--Wiles/Poitou--Tate formula therefore
gives
\[
 \frac{|\operatorname{Sel}_{L}(M_n)|}
      {|\operatorname{Sel}_{L^\perp}(M_n^\vee(1))|}
                       =|P_n|=|k|^n,
 \qquad |\operatorname{Sel}_{L}(M_n)|\ge|k|^n.              \tag{BE3}
\]
These are the actual finite Selmer groups with the stated
ordinary and unramified conditions.

Normalize a cocycle representative by \(z_n(c)=0\).
This representative is unique in its global cohomology
class: on \(M_n\), a global coboundary has value \(-2Y\)
at \(c\).
At \(\ell\), its correcting vector modulo \(P_n\) is
unique, since \(H^0((M/P)_n)=0\).
For the inverse limit use the exact inertia relations of the
deformation problem. Put
\[
 \mathcal N=\bigcup_{p\ne\ell}
                   \bigl(I_p\cap\ker r\bigr).
\]
Let \(Z_n\) be the set of continuous cocycles normalized by
\(z_n(c)=0\), satisfying \(z_n|_{\mathcal N}=0\), together
with the unique ordinary correcting-vector coset in
\(M_n/P_n\).
These sets are finite: every such cocycle is unramified
outside the fixed ramification set of \(r\), and
\(H^1(G_{\mathbb Q,S},M_n)\) is finite for a finite
coefficient module; normalization injects \(Z_n\) into
that finite group.
Every normalized unramified Selmer class counted in
(BE3) belongs to \(Z_n\), since an inertia coboundary
vanishes on \(I_p\cap\ker r\).
Consequently \(|Z_n|\ge|k|^n\).
The sets \(Z_n\) may be larger than that Selmer group;
only this lower bound is needed.

These sets have compatible primitive elements at every
precision.
If every value of a representative is divisible by \(\pi\),
divide it by \(\pi\) to obtain a unique representative at
precision \(n-1\).
The ordinary condition divides as well.
Indeed modulo \(\pi\) the correcting coset is invariant
in \(M/P\), so (BE1) makes it zero.
Subtract a vector of \(P\) from its lift and write the
remainder as \(\pi Y'\).
Dividing the local cocycle relation is then legitimate
because \(P\) is saturated.
The exact zero conditions on \(\mathcal N\) and at \(c\)
also divide.
Thus the nonprimitive part of \(Z_n\) injects into \(Z_{n-1}\).
If primitive representatives were absent at one precision,
they would be absent at all greater precisions, and this
injection would bound all subsequent \(|Z_n|\) by one fixed
finite cardinality, contradicting (BE3).
Every precision therefore has a primitive representative.
Their finite inverse-system tree has vertices at every
height; the finite branching argument gives a compatible
infinite branch.

Completeness now produces an actual nonzero continuous
cocycle \(z:G_{\mathbb Q}\to M\), with \(z(c)=0\), and
a correcting vector \(Y\in M\) such that
\[
                  u^{-1}(z-dY)u
                           \text{ is strictly upper locally}. \tag{BE4}
\]
The compatible correcting cosets lift because \(M/P\) is
finite free and complete.
The exact zero relations in the definition of \(Z_n\) give
\[
                 z(g)=0\quad(g\in I_p,\ r(g)=1),          \tag{BE5}
\]
for \(p\ne\ell\), at every precision and hence integrally.
We do not assert that arbitrary division preserves the
smaller unramified Selmer group; it is the exact
killed-inertia relation (BE5) that is needed and preserved.
Multiply \(z,Y\) by \(\pi\) once. They remain nonzero where
required and now reduce to zero.

Define
\[
                         r'(g)=(1+z(g))r(g),\qquad
                         u'=(1-Y)u .
\]
Because products of two elements of \(M\) are zero, the
cocycle equation is exactly the representation equation
for \(r'\).
The equations \(z^\#=-z\), \(Y^\#=-Y\) and square-zero
multiplication preserve the multiplier and make \(u'\)
symplectic.
Equation (BE4) gives an upper triangular local matrix:
\[
 (u')^{-1}r'(g)u'
    =u^{-1}(1+z(g)-dY(g))r(g)u .
\]
Its diagonal is unchanged, so all actual ordered inertial
characters are unchanged.
Equation (BE5) gives the required killed-inertia relations.
Reduction is unchanged because \(z,Y\in\pi M\).

The matrices \(r'\) are block triangular with the same
diagonal blocks as \(r\), hence have the same characteristic
polynomials.
The symplectic pairing identifies the two diagonal blocks
as duals up to the fixed multiplier; irreducibility of one
therefore implies that of the other.

Finally a splitting would give an off-block matrix \(Q\)
with \(z(g)=r(g)Qr(g)^{-1}-Q\).
At \(c\) its right-hand side is \(-2Q\), while \(z(c)=0\).
Thus \(Q=0\), forcing \(z=0\), a contradiction.
This proves the exact nonsplit \(\mathrm{BlockExt}\)
conclusion with all local and integral properties.
\end{proof}

The fixed-precision cohomological formula used above is
the same proved finite-module formula used in the arithmetic
Selmer estimates; neither automorphic adjoint vanishing
nor a lift of the original rank-two block was used here.
Finite residue fields, finite cocycle sets and the compatible
inverse limit give the required complete DVR model.
The source producers expanded are
\srcid{BlockMatrixPlaces.GwIneq}, \srcid{GwCount},
\srcid{PrimEx}, \srcid{DivStep}, \srcid{Compat},
\srcid{CocT}, and \srcid{CocA}, assembled in
\srcid{CompactIntegralFields.NotOver.bl}.
